\subsection{An exhaustive split of the affine-linear coordinate}\label{sec:wsplit}

Write $h=a+bZ+cT+dZ^2+eZT+fT^2$ and
$\Delta=e^2-4df$. The following split is exhaustive and its branches
are disjoint. The last column distinguishes a proof given here from
the proposition proving that branch. Every row is now proved here.

\begin{center}\small
\begin{tabular}{@{}p{0.27\textwidth}p{0.38\textwidth}p{0.25\textwidth}@{}}
\toprule
Coefficient condition & Reduction & Status in this draft\\
\midrule
$\Delta\ne0$ & Laurent phase on $\C^*$ with logarithmic volume &
Proposition~\ref{prop:nondeglog}\\[3pt]
$\Delta=0$, $f\ne0$ & $h=a+bZ+cT+fT^2$ after a null rotation &
Proposition~\ref{prop:rankoneformal}\\[3pt]
$f=e=0$, $d\ne0$, $c\ne0$ &
$W(a+T+Z^2)+G$ after normalization &
Proposition~\ref{prop:mixedformal}\\[3pt]
$f=e=0$, $d\ne0$, $c=0$ & $Wh(Z)+G$ &
Proposition~\ref{prop:height}\\[3pt]
$d=e=f=0$, $c\ne0$ & $W(a+bZ+cT)+G$ &
Corollary~\ref{cor:linearwt}\\[3pt]
$d=e=f=c=0$ & $Wh(Z)+G$ &
Proposition~\ref{prop:height}\\
\bottomrule
\end{tabular}
\end{center}

Here is the algebra behind the reductions. The null rotation
\[
 Z=Z',\quad T=T'+\lambda Z',\quad
 W=W'-\lambda T'-\lambda^2Z'/2
\]
preserves $2ZW+T^2$. If $f\ne0$, choose $\lambda=-e/(2f)$.
When $\Delta=0$ the quadratic part of $h$ becomes just $fT'^2$.
If $f=e=0$ and $c\ne0$, choose $\lambda=-b/c$ to remove the
linear $Z$ term without changing $dZ^2$. If also $d\ne0$, the
Gaussian scaling $Z=\kappa Z'$, $W=W'/\kappa$, followed by a
nonzero scaling of $P$, makes the $T$ and $Z^2$ coefficients both
one; take $\kappa^2=c/d$. All additional terms enter a polynomial
$G$ of degree at most three. If $\Delta=0$ and $f=0$, then $e=0$,
which proves that the table omits no case.

\subsection{A curve-length estimate}

\begin{lemma}\label{lem:curvelength}
A real algebraic curve of degree at most $D$ has length at most
$\pi D\rho$ in a disk of radius $\rho$. Consequently its logarithmic
length in $a\leq|x|\leq b$ is at most a constant times
$D(1+\log(b/a))$.
\end{lemma}
\begin{proof}
Average the number of intersections with lines over directions and
apply the one-dimensional coarea formula. Outside exceptional line
components the number is at most $D$; the components themselves
satisfy the same estimate. In each dyadic annulus
$\rho\leq|x|\leq2\rho$, divide the disk estimate by $\rho$ to
bound $\int|dx|/|x|$ by $2\pi D$. Sum the dyadic annuli.
\end{proof}

\subsection{The logarithmic branch, including its removed component}

For $\Delta\ne0$, $f\ne0$, choose $\eta^2=\Delta$ and set
\[
 \rho=\frac{e+\eta}{2f},\quad \sigma=\frac{e-\eta}{2f},\quad
 u=\frac{\rho c-b+1/t}{\eta},\quad
 v=\frac{b-1/t-\sigma c}{\eta},\quad k=uv-a/f.
\]
On the circular constraint $z=th(z,T)$, put
$x=T+\rho z+u$, so that $T+\sigma z+v=k/x$. Then
\[
 z=\frac f\eta(x-k/x-u+v),\quad T=x-u-\rho z,
 \qquad
 \frac{dT}{1-th_Z}=\frac{dx}{t\eta x}.
\]
The coefficient of $t^{-2}$ in $k$ is $-1/\eta^2$, so the conic
is generically smooth. The exponent $tG(z,T)-T^2/2$ is Laurent
polynomial in $x$. On the formal Gaussian branch,
$x=1/(\eta t)+O(1)$, which makes every Laurent power legitimate.

If $f=0$, then $e\ne0$. Put $z_0=-c/e$ and $x=z-z_0$. With
$A(z)=a+bz+dz^2$, the constraint gives
\[
 T=\frac{z/t-A(z)}{ex},\qquad
 \frac{dT}{1-th_Z}=\frac{dx}{tex}.
\]
This localization is valid if $c\ne0$, or if $c=0$ and $a\ne0$.
If $a=c=0$, the chosen fixed point is instead $g=0$ identically;
inverting $g$ would discard the very branch being tested. In that
case the exact unit series is
\[
 \E_T\left[\frac{e^{tG(0,T)}}{1-t(b+eT)}\right].
\]
Use $x=T-(1/t-b)/e$, now a valid Laurent unit. Its volume is
$-dx/(tex)$ and its exponent is polynomial in $x$. These formulas
cover both the smooth conic and the removed-component branch.
We next connect these coordinates to the formal Gaussian functional.
This will close every one-ended phase without an analytic estimate.

\begin{lemma}[Formal logarithmic transport]\label{lem:logtransport}
In each of the three charts above, write $X=X(t,T)$ for the formal
coordinate, $D$ for the displayed circular denominator, and $B$ for
$tG(g,T)$ or $tG(0,T)$ on the removed component. Put
$\kappa=\eta,e,-e$, respectively. Then
\[
 J(q)=\E_T\!\left[\frac{q(X)e^B}{D}\right]
\]
is horizontal on the logarithmic twisted quotient with phase
$\varphi$ and normalization $h=1/t$, and $J(1)=\Phi_1$.
\end{lemma}
\begin{proof}
All Laurent powers of $X$ are legitimate by the unit checks above.
The coordinate identities give
\[
 X_T=t\kappa X/D,\qquad
 \varphi(t,X)=B-T^2/2.
\]
Consequently Gaussian integration by parts gives
$J((x\partial_x+x\varphi_x)q)=0$: the corresponding integrand,
multiplied by $t\kappa$, is
$(\partial_T-T)\{q(X)e^B\}$.

For parameter horizontality set $w=1/D$ and $R=X_t/X_T$.
Differentiating $w=X_T/(t\kappa X)$ gives the exact identity
\[
 w_t+w/t=\partial_T(wR).
\]
Subtracting $J(q_t+(\varphi_t-1/t)q)$ from $\partial_tJ(q)$ now
gives the Gaussian expectation of
$(\partial_T-T)\{wR q(X)e^B\}$, hence zero. These are formal
coefficientwise identities, without an integral-convergence hypothesis.
The seed equality is the circular contraction formula.
\end{proof}

\begin{proposition}\label{prop:oneendgaussian}
Every chart above whose actual Laurent phase has at most one
irregular end is not pure. The removed-component branch
$a=c=f=0$, $e\ne0$, is therefore excluded for every polynomial $G$,
without a degree bound.
\end{proposition}
\begin{proof}
At most one irregular end means that $\varphi$ is polynomial in
$x$ or in $x^{-1}$, including the constant case. Purity would give
the horizontal functional in Lemma~\ref{lem:logtransport} seed value
one, contrary to Corollary~\ref{cor:oneend}. The removed-component
phase is always polynomial in $x$, irrespective of $\deg G$.
\end{proof}

\subsection{Resolving a Laurent phase}\label{sec:laurentresolution}

The logarithmic estimate needs two ingredients: a complete description
of the regular part, and a bound for the distinguished differential
$d\log x$. The first is a nearby-cycle calculation. The second is a
calculation on annuli. Keeping these ingredients separate prevents a
perfect period pairing from being mistaken for a growth estimate.

Throughout this subsection the coefficients are rational functions of
the parameter. In particular their Puiseux expansions converge. The
nonconvergent formal statements proved earlier have a different scope.

The two-ended twisted quotient is finite dimensional without a
genericity assumption. If the actual Laurent orders are $p,q\geq1$,
the classes $x^{-q},\ldots,x^{p-1}$ are a basis. The top term of
$x\varphi_x$ eliminates every exponent at least $p$, without crossing
below $-q$; the bottom term eliminates every exponent below $-q$,
without crossing above $p-1$. An exact expression supported in this
interval would require its preimage to have highest exponent at most
$-1$ and lowest exponent at least zero. This proves independence
and rank $p+q$. Coefficient specializations use their actual orders.

\begin{lemma}[A marked model]\label{lem:markedmodel}
Let $\varphi(r,x)\in\C(r)[x,x^{-1}]$ be nonconstant in $x$.
After a finite ramification of $r$ there is a smooth surface proper
over a parameter disc, unchanged over its punctured disc, with these
properties:
\begin{enumerate}
\item Its generic fibre is $\mathbb P^1$. Its central fibre is a
reduced normal-crossing curve, with local equations $r=u$ or $r=uv$.
\item The phase extends to a morphism to $\mathbb P^1$. Its zero
and pole divisor and the divisor of $x$ have normal crossings with
the central fibre.
\item The horizontal sections $x=0,\infty$ are disjoint and meet
smooth points of the central fibre. Locally $x=u^a v^b U$, where
$a,b$ are integers and $U$ is holomorphic and nonvanishing.
\end{enumerate}
Additional markings used in the construction are not additional
punctures of the exponential connection.
\end{lemma}
\begin{proof}
Mark $0,\infty$ and every point of the generic fibres
$\varphi^{-1}(0),\varphi^{-1}(1),\varphi^{-1}(\infty)$.
Retain multiplicities in the divisors, but identify identical sections.
Convergent Newton--Puiseux makes the finitely many marked branches
meromorphic sections after one finite ramification. Start with
$\mathbb P^1$ over that disc and blow up collisions of distinct
sections. This terminates: the contact order of each pair is a finite
integer and decreases on blowing up their common point.

Every centre is a smooth point of a single central component. Indeed
a section cannot pass through a node $uv=r$: two positive integral
orders cannot have sum one. Thus each new exceptional component has
multiplicity one and all nodes still have the equation $uv=r$.
The final marked sections are disjoint and miss the nodes.

The divisor of the meromorphic function $\varphi$ is supported on
the central components and its marked zero and pole sections. If a
point were an indeterminacy, its local divisor would have a zero
branch and a pole branch meeting there. After cancelling common
factors its expression is a unit times $u^a/v^b$, $a,b>0$.
The equation $\varphi=1$ then has a horizontal branch through that
point. At a node this contradicts the section property; at a smooth
point it contradicts the separation from the zero or pole section.
These are all possible intersections. Hence $\varphi$ is a morphism.
The asserted normal crossings and the expression for $x$ follow
from its divisor: horizontally $x$ has just the two marked sections.
Only those original sections are removed from the generic fibre;
marking a zero or a point with phase value one does not remove it.
\end{proof}

The marking of phase value one is necessary. For $\varphi=r/x^2$,
the zero and pole sections are already disjoint, but the map is
undefined at $r=x=0$. Its value-one branches $x=\pm\sqrt r$
force the cover $r=s^2$ and the chart $x=sy$, where the phase is
$1/y^2$. The marked construction therefore detects a transition
that merely separating zeros and poles would miss.


We use the following precise comparison facts. For a meromorphic
connection on a punctured disc, its moderate nearby cycles describe
its formal regular summand, with monodromy. Nearby cycles commute
with proper direct image. For a rank-one exponential on a surface
with normal crossings, a component on which the phase has a vertical
pole contributes no moderate nearby cycles. At a node between a
finite-phase component and a pole component, the surviving object is
the exponential connection on the finite component, localized at the
node; it is not an extra point class. At a node between two finite
components, the local complex is the usual nearby annulus complex,
including its monodromy. These are standard comparison and
normal-crossing facts, not new theorems here; see
Roucairol~\cite[Section~1.2 and Lemma~1.4]{RoucairolFormal} and
Sabbah~\cite[Section~1.b, especially Proposition~1.5]{SabbahStationary}.

For clarity, their local exponential models are a holomorphic unit
(finite phase), $\exp(u^{-p})$ at a finite--pole boundary, and
$\exp(u^{-p}v^{-q})$, $p,q>0$, when both components are poles.
An analytic unit in a pole monomial is absorbed in a coordinate.
This may multiply the equation of the base by a nonvanishing unit;
locally that change does not change the vanishing assertion or the
nearby-cycle object up to the usual unit identification. In the
$f^s$ description the identification is multiplication by
$\exp(s\log U)$, with a local analytic choice of $\log U$.
One keeps the actual gluing maps and monodromy, rather than just
adding the ranks of the components.

\begin{lemma}[Lifting the regular part]\label{lem:regularlifting}
For the marked model, the complete nearby-cycle space of the
exponential direct image is represented, on a sufficiently small
parameter sector, by flat rapid-decay chains made from the following
pieces:
\begin{enumerate}
\item bounded deformations of finitely many paths and cycles in
the finite-phase central components, away from their nodes and poles;
\item a bounded number of bridges and circles in the necks
$uv=r$ between finite-phase components;
\item descending phase tails at irregular ends and circles about
regular punctures.
\end{enumerate}
The coefficients of each chosen flat chain are fixed constants.
Their nearby-cycle classes exhaust the formal regular part, with
its actual gluing maps and monodromy.
\end{lemma}
\begin{proof}
Apply proper nearby-cycle comparison to the compactified surface,
with the exponential connection localized at $r=0$ and at the two
original punctures. The normal-crossing rules just stated kill the
interiors of the pole components. Along a finite component they leave
its ordinary exponential connection, with localization at its pole
ends and at any original regular puncture. A finite--pole node gives
that localized end. A finite--finite node retains the nearby annulus.
Thus no pole component is silently counted as an additional regular
class, and no annulus or puncture residue is discarded.

Here is a chain description of this calculation. Remove small discs
about the finitely many nodes and punctures in each finite component.
Triangulate the remaining compact surfaces, keeping the boundary
circles. At an irregular boundary attach the finitely many descending
sectors. At a regular puncture keep the circle. The one-dimensional
exponential period comparison identifies the resulting relative
chain complexes with the component de Rham complexes; for a constant
phase this is ordinary homology. At each finite--finite node glue by
the actual annulus chain complex, with both its degree-zero and
degree-one terms. Taking the total complex is Mayer--Vietoris for
the full nearby-cycle complex above. It computes its maps and
monodromy, not merely its Euler characteristic. In particular its
classes exhaust the formal regular summand by proper comparison.

Choose a finite basis of this finite-dimensional homology. Smooth
pieces and their boundary points lift by the implicit function
theorem; an annulus is lifted by $uv=r$. Descending ends are continued
by the inverse phase coordinate. On a fixed small sector these choices
are isotopies, so the homology coefficients are constants. No Morse
condition on the phase is used. Degenerate critical points stay
inside the triangulated pieces, and a constant phase on a component
is allowed.

For the tail continuation, choose the ray of $r$ so that no unbounded
critical value lies on the negative real phase axis for small $r$.
There are finitely many excluded directions: critical values are
convergent Puiseux functions. Bounded critical values fit in a fixed
disc. From a sufficiently negative phase value every chosen inverse
branch therefore continues to $-\infty$ without meeting a critical
value. In the two-ended Laurent case it ends at $x=0$ or $\infty$.
The same choice can avoid the finitely many limiting Stokes directions
of the parameter connection.

This constructs flat lifts whose nearby-cycle classes exhaust the
space computed by proper comparison. No quantitative estimate has
been used in that completeness calculation.
\end{proof}

\begin{lemma}[The logarithmic bound on the lifted pieces]
\label{lem:logpieces}
For the chains of Lemma~\ref{lem:regularlifting}, the period of
$e^{\varphi}d\log x$ is $O(1+|\log r|)$ along the chosen ray.
The bound is uniform for a fixed finite collection of generators.
With a fixed Laurent amplitude the same construction has at most
an additional algebraic power of $|r|^{-1}$.
\end{lemma}
\begin{proof}
On a compact finite-phase piece the phase and the coefficients of
$d\log x$ are bounded, so its integral is bounded. In a finite--finite
neck write
\[
 r=uv,\qquad x=u^a v^b U(u,v).
\]
At fixed $r$,
\[
 d\log x=\left(a-b+\frac{uU_u-vU_v}{U}\right)\frac{du}{u}.
\]
The coefficient is uniformly bounded in a small bidisc. A radial
bridge across $|r|/\eta\leq|u|\leq\eta$, with bounded angular
detours, has logarithmic length $O(1+|\log r|)$. Circles have bounded
logarithmic length. The phase is holomorphic on this bidisc and is
uniformly bounded there. This treats every finite--finite neck.

Near an irregular end of a finite component, use a local root of
the phase unit to write $\varphi=w^{-p}$, $p\geq1$. At a
finite--pole node one can simultaneously write $r=wv$ after changing
the other coordinate by a unit. On $|v|,|w|<\eta$, the same
calculation bounds $d\log x$ by a constant times $|dw|/|w|$.
Along a descending phase ray, $s=-\varphi>0$, its integral is bounded
by a constant times
\[
 \int_{s_0}^{\infty}e^{-s}\frac{ds}{s}.
\]
If the node chart ends, truncate it first at $|w|=|r|^\alpha$,
where $0<\alpha<1$. Its phase is already of size
$-|r|^{-\alpha p}$. Continue that inverse branch along the negative
phase axis as in Lemma~\ref{lem:regularlifting}.

We give a uniform bound for that continuation, rather than declaring
it negligible from a local saddle formula. Choose $M$ large enough
to include all coefficient transition scales in
$|r|^M\leq|x|\leq|r|^{-M}$. Such $M$ exists because the coefficients
are rational and there are finitely many Laurent powers. In this
annulus each descending path lies on $\operatorname{Im}\varphi=0$,
a real algebraic curve of bounded degree after clearing its Laurent
denominator. The curve-length bound used in
Lemma~\ref{lem:curvelength} gives
$\int|dx|/|x|=O(1+|\log r|)$, by dyadic annuli. Along the entire
tail the exponential is at most
$\exp(-|r|^{-\alpha p})$. Outside that annulus the leading Laurent
monomial dominates uniformly, so
$|d\log x/d\varphi|\leq C/|\varphi|$ and its descending integral
is bounded by the same exponential tail integral. This proves the
required uniform tail bound. A fixed Laurent amplitude contributes
only a power of $|r|^{-1}$ on the annulus and algebraic growth in
$s$ outside; the exponential still dominates it.

At a regular puncture use a small fixed circle in its smooth marked
chart. Both its phase and logarithmic length are bounded. There are
finitely many pieces and their chain coefficients are fixed.
Adding their bounds proves the lemma.
\end{proof}

\begin{theorem}[Logarithmic seed bound]\label{thm:laurentlog}
Let $\varphi(t,x)\in\C(t)[x,x^{-1}]$. On a generic ray at
$t=\infty$, the zero-exponential part of every scalar flat period
of $e^{\varphi(t,x)}dx/x$ is $O(1+\log|t|)$. This includes all
critical degeneracies, coincident actions and regular punctures.
Consequently, for a rational normalization $h(t)$ with
$v_\infty(h)>0$, its differential seed submodule admits no rational
horizontal quotient taking the seed to one. The same exclusion
holds for a horizontal formal Gaussian functional at another point.
\end{theorem}
\begin{proof}
If the phase is independent of $x$, the sole residue period is a
constant times $e^{\varphi(t)}$. Its zero-exponential part is bounded
when $\varphi$ is regular at infinity, and absent otherwise.
For a nonconstant phase put $r=1/t$ and use
Lemmas~\ref{lem:markedmodel}--\ref{lem:logpieces}. The construction
includes regular puncture circles when there is just one irregular end.
Finite ramification changes $\log|r|$ only by a constant factor.
The lifted chains have at most algebraic growth on every Laurent
amplitude by Lemma~\ref{lem:logpieces}. On the chosen generic ray,
a full period vector with such growth has no exponentially growing
factor: after formal exponential decomposition a largest growing
factor could not cancel on every basis amplitude. The other nonzero
factors are exponentially decaying there. The nearby-cycle basis
therefore gives a full basis of regular functionals modulo those
decaying terms, with constant coefficients in the flat basis. Every
one of their seed periods has the asserted logarithmic bound. Thus
their zero-exponential seed
expansions have no negative algebraic powers. Multiplication by $h$
makes all of them tend to zero.

For a constant phase the rank-one period already gives the same
contradiction directly. For a nonconstant phase, a quotient with
normalized seed one would extend locally to a
horizontal functional on the full module, by the seed restriction
and analytic extension argument of Lemma~\ref{lem:seedfunctional}.
Its comparison hypotheses hold on
$(\mathbb A^1\setminus S)\times\mathbb G_m$: remove the finitely
many parameter poles, zeros of $h$, and comparison exceptions.
The base projection and nonconstant phase are regular and
algebraically independent there.
Its scalar seed is constant. Distinct formal exponential factors
separate, so its zero-exponential seed is also one. This contradicts
the bound after normalization. A formal Gaussian functional with
unit seed first restricts to the same rational seed quotient, as in
the first part of Lemma~\ref{lem:seedfunctional}.
\end{proof}

The comparison theorems supply completeness here. The bounds come
from the logarithmic differential on the marked model and the explicit
descending-tail estimate. Neither part follows from a finite symbolic
calculation. This proof uses convergence of the rational parameter
family; it makes no assertion for arbitrary nonconvergent Laurent
families or arbitrary coefficient rings.


\begin{proposition}\label{prop:nondeglog}
If $e^2-4df\ne0$, no polynomial $G$, of any degree, makes
$P=Wh(Z,T)+G(Z,T)$ pure.
\end{proposition}
\begin{proof}
The smooth-conic, linear-conic and removed-component charts are
exhaustive. Their Laurent-unit checks and
Lemma~\ref{lem:logtransport} give the formal Gaussian functional
with seed value $\Phi_1$ and normalization $1/t$.
Theorem~\ref{thm:laurentlog} excludes $\Phi_1=1$ in every chart,
using its actual Laurent orders. No image-kernel calculation is used.
\end{proof}

\begin{proof}[Proof of Proposition~\ref{prop:wlinear}]
Use the disjoint case table in Section~\ref{sec:wsplit}.
Proposition~\ref{prop:nondeglog} excludes its first row.
Proposition~\ref{prop:rankoneformal} treats the second and leaves
only a one-sign product. Proposition~\ref{prop:mixedformal} excludes
the third. Proposition~\ref{prop:height} proves the fourth and sixth;
its pure cases are one-sign phases or the shear family.
Corollary~\ref{cor:linearwt} treats the fifth and leaves a quadratic.
These rows exhaust all coefficients, including their zero faces.

The one-sign and ternary quadratic cases have cutoff $m>\deg Q$,
and the shear has cutoff $m>3\deg Q$. The null rotations,
Gaussian coordinate scalings and nonzero scalar multiples of $P$
used in the reductions preserve the threshold and multiplier degree.
Thus every pure member of the full affine-linear family satisfies
$m>3\deg Q$.
\end{proof}

\subsection{Adversarial tests for the analytic argument}

The following exact examples prevent several tempting shortcuts.
For positive real $t$,
\[
 \frac1t\int_{\mathbb R}e^{-x^2/(2t^k)}\,dx
   =\sqrt{2\pi}\,t^{k/2-1}.
\]
The unique critical point is bounded and has zero action for every
$k$, yet the normalized period decays for $k=1$, is constant for
$k=2$, and grows for $k=3$. A critical-point count or a zero-action
label alone gives no decay estimate.

For the fixed phase $tx^k$, $k>2$, the critical point is degenerate
and the decay scale is $t^{-1/k}$, not a Gaussian Hessian scale.
The fixed quotient in Theorem~\ref{thm:stationaryballs}
retains its full multiplicity $k-1$.
For $tV_nx^n-x^4/2$, $1\leq n\leq3$, the $n-1$ bounded and
$4-n$ escaping critical dimensions total three. Counting only
bounded points loses actual rank; the stationary-ball construction
retains all roots and tests its full matrix using polynomial separation.
Finally, $t^{-3/2}=r^3$ under $r=t^{-1/2}$: a claim about a
noninteger exponent before ramification must not be reused after
ramification. Lemma~\ref{lem:positiveformal} uses positive seed order
and an actual full frame instead.

\subsection{One exclusion criterion for several geometric cases}
\label{sec:gradientexclusion}

The smooth case does not need a list of critical-point formulas.
The same argument also treats many singular leading forms. Its
hypotheses must hold in every complex direction. A bound only on
the original real integration space would not suffice.

We state the criterion for a polynomial exponential amplitude, then
specialize it to the Gaussian. Norms are Hermitian norms on complex
affine space; the Gaussian quadratic remains holomorphic.

\begin{theorem}[A gradient criterion]\label{thm:gradientcriterion}
Let $P,q\in\C[x_1,\ldots,x_n]$, with $P$ nonconstant and
$\deg q\leq J$, where $J\geq1$. Suppose that outside a compact set
\begin{equation}\label{eq:gradientconditions}
 \|\nabla P(x)\|\geq c\|x\|^{J-1},\qquad
 \frac{\|\operatorname{Hess}P(x)\|}{\|\nabla P(x)\|^2}\leq C
\end{equation}
for some $c>0$ and $C<\infty$. For every fixed polynomial $A$,
no nonzero flat period of
\[
 A(x)e^{tP(x)-q(x)}\,dx_1\cdots dx_n
\]
is a finite exponential-polynomial sum
$\sum_\lambda e^{\lambda t}p_\lambda(t)$ with polynomial
$p_\lambda$. In particular it cannot be a nonzero constant.
\end{theorem}
\begin{proof}
We first justify the full cycle space and the tails. The gradient is
nonzero outside a compact set, so the affine critical set is compact
and algebraic, hence finite. The gradient bound is stronger than the
condition excluding vanishing cycles at infinity. The finite-critical
thimble decomposition therefore applies, with the whole local Milnor
homology at each isolated critical point, not just its Morse
approximation~\cite[Part~I, Sections~1--4]{Pham}.
When several critical points have the same value, take disjoint
Milnor balls and apply excision to their complement. This gives the
direct sum of their full local homology before any values are grouped.

Here the Gaussian or polynomial amplitude can be kept on these fixed
$P$-cycles. Choose a generic argument $\theta$ of $t$. Outside the
local critical neighbourhoods, lift a descending phase ray by
\[
 P(x)=\lambda-e^{-i\theta}s,\qquad
 \frac{dx}{ds}=-e^{-i\theta}
                    \frac{\overline{\nabla P(x)}}{\|\nabla P(x)\|^2}.
\]
For large $s$ the lift is outside the compact set, because its phase
value is unbounded. The first condition in
\eqref{eq:gradientconditions} gives
$\|x\|^J\leq C_1(1+s)$. The real differential of this vector field
has norm at most a constant times the second ratio in
\eqref{eq:gradientconditions}, hence is bounded. Transported
$(n-1)$-cycle volumes grow at most as $\exp(C_2s)$; the radial speed
is bounded as well. Also
$|A(x)e^{-q(x)}|\leq C_3(1+s)^M\exp(C_4s)$ for some $M$.
Thus $\exp(-|t|s)$ dominates the complete tail, uniformly for large
$|t|$ in a smaller sector. The same estimate applies to every fixed
polynomial amplitude.

We must also check that these are a full basis for the actual phase,
rather than selected admissible cycles. Put
\[
 F_\sigma=P-\sigma q/t,\qquad 0\leq\sigma\leq1.
\]
Since $\|\nabla q\|\leq C_5\|\nabla P\|$ outside the same compact
set, all large $|t|$ give
\[
 \|\nabla F_\sigma\|\geq\tfrac12\|\nabla P\|.
\]
The ratio involving the Hessian stays uniformly bounded too. Indeed
$\|\operatorname{Hess}q\|=O(\|x\|^{J-2})$ for $J\geq2$, while
that Hessian is zero for $J=1$.
All critical points and their values therefore remain in fixed compact
sets. On the preimage of a phase half-plane beyond those values, the
parameter field
\[
 W_\sigma=\frac{q}{t}
       \frac{\overline{\nabla F_\sigma}}{\|\nabla F_\sigma\|^2}
\]
satisfies $\partial_\sigma F_\sigma+dF_\sigma(W_\sigma)=0$ and
$\|W_\sigma\|\leq C_6(1+\|x\|)/|t|$. Its flow is complete for
$0\leq\sigma\leq1$. Extend it by a cutoff away from that half-plane.
It gives an isotopy of the ambient relative pairs defining the full
rapid-decay homology. This explicitly prevents an infinity class
from appearing during the deformation.

The fixed $P$-thimbles have uniformly descending tails for every
$F_\sigma$, by the preceding estimates. Truncate them beyond a common
large phase value. Their boundaries stay in the descending half-plane
throughout the isotopy, so they define its flat relative classes.
They are a full basis at $\sigma=0$ and remain one at $\sigma=1$.
Finite critical points may merge, and their critical values may
coincide; the argument uses the full relative pairs and does not
require a Morse family. The same cycles are flat as $t$ varies in
the chosen sector.

These improper chains and their homotopies are admissible for the
period pairing. One can integrate first on truncated chains and then
pass to the limit: the tail bound makes every boundary term tend to
zero. Indeed the differential of $W_\sigma$ grows at most as
$C(1+\|x\|^J)/|t|$; along the controlled tails this is
$C'(1+s)/|t|$. The homotopy volumes have at most this exponential
growth, still dominated by the descending factor for all large $|t|$.
Thus Stokes and differentiation under the integral hold throughout;
the relative-pair calculation does not depend on a convergence claim
for the original real Gaussian integral.

It remains to determine the local asymptotics at a finite critical
point with value $\lambda$. For an analytic amplitude germ, the
vanishing-cycle density is a finite sum of terms
$s^{\alpha+k}(\log s)^j$, with analytic coefficient series and
$\alpha>-1$. Its Laplace asymptotics consequently have the form
\[
 e^{\lambda t}\sum_{\alpha,k,j}
       c_{\alpha,k,j}\,t^{-\alpha-k-1}(\log t)^j,
 \qquad k\geq0.
\]
This is the local analytic result of
Pham~\cite[Part~II, Section~1.1, (1.5)--(1.6) and the footnote on
printed page~26]{Pham}; it permits degenerate isolated critical
points. Pham credits Maisonobe and Rombaldi for the local solution
comparison used there. The amplitude $A e^{-q}$ is analytic. Every nonzero
term has a strictly negative power of $t$ and tends to zero after
removing $e^{\lambda t}$. The tails outside the local neighbourhood
are exponentially negligible by the uniform bound already proved.

The finite-critical decomposition accounts for the full module.
Group equal critical values before separating the distinct formal
exponential factors. Within each such factor the displayed local
expansions contain only decaying powers and logarithms. A nonzero
polynomial prefactor cannot occur. The zero function is allowed when
all these contributions vanish. This proves the theorem.
\end{proof}

\begin{corollary}[Original Gaussian markers]\label{cor:gradientmarkers}
For the standard Gaussian quadratic $q=\sum x_i^2/2$, suppose
\eqref{eq:gradientconditions} holds with $J=2$. Then every original
marked series $\Phi_A$ is either zero or is not a finite
exponential-polynomial sum. In particular $P$ is neither pure nor safe.
\end{corollary}
\begin{proof}
The exact top de Rham quotient and connection are
\[
 M=\C(t)[x]\Big/\sum_i(\partial_i+tP_i-x_i)\C(t)[x],
 \qquad\nabla[A]=[A_t+PA].
\]
Gaussian integration by parts makes $[A]\mapsto\Phi_A$ horizontal
on this vector-space quotient. The spatial and parameter derivatives
commute; no ordinary Jacobian-ideal quotient is substituted.
If some $\Phi_A$ is a
nonzero finite exponential-polynomial sum, a constant-coefficient
differential operator isolates one exponential and kills its
polynomial prefactor, giving $c e^{\beta t}$, $c\ne0$.
Because
\[
 \partial_t\Phi_A=\Phi_{PA},
\]
this is still an original fixed polynomial marker. Replace $P$ by
$P-\beta$ and divide the marker by $c$. Its formal seed value is one.
Lemma~\ref{lem:seedfunctional} gives an actual constant period for
that marker, contradicting Theorem~\ref{thm:gradientcriterion}.
Its comparison hypotheses hold with base $t$ and phase $tP-q$ on
$\mathbb A^1\times\mathbb A^n$, after removing the finite exceptional
base set. No convergence of the original Gaussian series is assumed.
Finally $\Phi_1(0)=1$, so it is a nonzero series; purity or safety
would make it a polynomial.
\end{proof}

\begin{corollary}[Smooth leading form, in every degree]
\label{cor:smoothleading}
If $\deg P=D\geq2$ and the homogeneous leading form satisfies
$\nabla F_D(v)=0$ only for $v=0$, then $P$ is not pure. In fact
Corollary~\ref{cor:gradientmarkers} holds for every original marker.
\end{corollary}
\begin{proof}
Compactness of the complex unit sphere gives
$\|\nabla P(x)\|\geq c\|x\|^{D-1}$ for large $\|x\|$.
Also $\|\operatorname{Hess}P\|=O(\|x\|^{D-2})$, so its ratio to
the squared gradient is $O(\|x\|^{-D})$. Both conditions for $J=2$
follow, with all lower terms retained.
\end{proof}

\begin{lemma}[The ordinary-node test]\label{lem:nodecoercivity}
Let $P=F_3+F_2+F_1+F_0$ in $n\geq2$ variables. Suppose its
projective singular set is finite and every singular point $[v]$
has $\ker\operatorname{Hess}F_3(v)=\C v$. If $F_2(v)\ne0$ at
every such point, then
$\|\nabla P(x)\|\geq c\|x\|$ outside a compact set. Consequently
$P$ is not pure under the Gaussian functional.
\end{lemma}
\begin{proof}
Otherwise choose an unbounded sequence with
$\|\nabla P(x)\|/\|x\|\longrightarrow0$. A projective limit $[v]$
is singular for $F_3$. Choose a fixed complement to $\C v$ and write
$x=sv+y$, where $y=o(s)$. With $H=\operatorname{Hess}F_3(v)$,
homogeneity gives the exact expansion
\[
 \nabla P(sv+y)=sHy+\nabla F_3(y)+s\nabla F_2(v)
                              +\nabla F_2(y)+\nabla F_1.
\]
The map $H$ is injective on the chosen complement. Divide by $s$.
The error from $\nabla F_3(y)/s$ is $o(\|y\|)$, so the equation
first forces $y$ to stay bounded. A limit $y_0$ satisfies
$Hy_0+\nabla F_2(v)=0$. Pair bilinearly with $v$. Symmetry and
$Hv=0$ kill the first term, whereas Euler's identity gives
$v\cdot\nabla F_2(v)=2F_2(v)$, a contradiction.
The Hessian of a cubic is $O(1+\|x\|)$, so the second condition
in \eqref{eq:gradientconditions} follows from the first. Apply
Corollary~\ref{cor:gradientmarkers}.
\end{proof}

\subsection{A common residual family for ordinary nodes}
\label{sec:ordinaryresidual}

For a pure cubic with only ordinary leading nodes, the preceding
lemma forces $F_2$ to vanish at some node. An isotropic such node
gives the frame already proved in Lemma~\ref{lem:frame}. Every
nonisotropic such node has the following single normal form.

\begin{lemma}\label{lem:ordinarynormal}
Let $[w]$ be an ordinary node of a ternary $F_3$, with $F_2(w)=0$
and $w^Tw\ne0$. An invertible affine change of coordinates puts
the pair consisting of $P$ and the Gaussian quadratic into
\begin{align*}
 P={}&T(uv+h)+a u^3+d v^3+b u^2+c v^2+e u+f v+k,\\
 q={}&\tfrac12(Au^2+Bv^2+mT^2+2Cuv+2DuT+2EvT)
                  +\ell_u u+\ell_v v+\ell_T T+q_0,
\end{align*}
where $m\ne0$ and $\det\operatorname{Hess}q\ne0$.
All Gaussian means and cross terms must be retained. Its leading
curve is irreducible nodal when $ad\ne0$, a transverse line and
smooth conic when exactly one of $a,d$ vanishes, and three
nonconcurrent lines when $a=d=0$.
\end{lemma}
\begin{proof}
Take $T$ in the direction $w$ and choose a complementary plane.
The singular-point equations make $F_3$ linear in $T$; its
coefficient is a nondegenerate binary quadratic because the node
is ordinary. Factor that quadratic and choose its two linear
factors as coordinates $u,v$. Thus $F_3=Tuv+B_3(u,v)$.
Replace $T$ by $T$ plus a linear form in $u,v$ to remove the
$u^2v$ and $uv^2$ terms of $B_3$.

There is no $T^2$ term in $F_2$, precisely because $F_2(w)=0$.
The full coefficient of $T$ is therefore
$uv+\alpha u+\beta v+h_0$. Translations of $u,v$ make it $uv+h$.
The cubic terms outside $T$ remain pure cubes. A constant shift
of $T$ removes the remaining quadratic $uv$ term. This gives the
displayed $P$. The direction of $T$ has not changed, so $m\ne0$;
invertibility of the affine change preserves the nondegenerate
Gaussian Hessian, while translations and the shifts of $T$ produce
the means and cross terms shown.

If $ad\ne0$, the polynomial $Tuv+a u^3+d v^3$ is irreducible:
an independent factor of a polynomial linear in $T$ would divide
both $uv$ and $a u^3+d v^3$, whose gcd is one. Its gradient shows
that its sole singular point is the stated ordinary node. If
$a=0,d\ne0$, it is $v(Tu+d v^2)$, with a smooth conic meeting
the line transversely at two points; the other one-zero case is
symmetric. If both vanish, it is the coordinate triangle $Tuv$.
\end{proof}

This consolidates three leading-curve rows into one residual family.
It is a reduction, not yet an exclusion of that family. In particular
the full quadratic $q$, including its means, cannot be replaced by
a convenient diagonal form during the remaining period argument.

\subsection{Two controls on the scope of the criterion}

Ordinary tameness of $P$ alone is insufficient for the Gaussian
conclusion. Let $P=\ell^Tx$, with $\ell=(1,i)/\sqrt2$ in two
variables. Its gradient is constant and nonzero, but
$\Phi_1=\exp(t^2\ell^T\ell/2)=1$. The phase $tP-q$ has its
critical point at $x=t\ell$, escaping with action zero. Thus an
isolated critical set for $P$, or the weaker gradient bound away
from zero, would lose exactly the contribution that purity needs.
Condition~\eqref{eq:gradientconditions} with $J=2$ prevents this.

The distinction between an original marker and a time-dependent
amplitude also remains essential. For
$P=(X_1^2+X_2^2)/2$ one has $\Phi_1=(1-t)^{-1}$ and
$(1-t)\Phi_1=1$. This phase satisfies the gradient criterion.
The unit in its rational differential module is represented by a
time-dependent amplitude. No original fixed polynomial marker has
a nonzero polynomial period, by Corollary~\ref{cor:gradientmarkers}.
The abstract module alone does not record that distinction.

\subsection{Ordinary nodes with a finite affine critical set}
\label{sec:ordinaryfinite}

The three ordinary-node families admit one argument whenever the
critical set of $P$ is finite. In particular we need not solve for
each escaping critical point, assume it is Morse, or separate its
action from the actions of other escaping points. The Gaussian
quadratic itself detects the unwanted contribution at infinity.

Retain the notation and the full pair $(P,q)$ of
Lemma~\ref{lem:ordinarynormal}. Write $q_2$ for the homogeneous
quadratic part of $q$. The absence of a Gaussian frame is exactly
the following condition at the leading nodes:
\begin{equation}\label{eq:nodenoframe}
 F_2(w)=0\quad\Longrightarrow\quad q_2(w)\ne0.
\end{equation}
Indeed the Gaussian bilinear form in these affine coordinates is
$2q_2$; the value of $F_2$ at a leading singular direction is
unchanged by translation. In the common normal form the $T$ node
already satisfies this condition because $m\ne0$. The only extra
conditions are $c\ne0$ or $B\ne0$ when $d=0$, and $b\ne0$ or
$A\ne0$ when $a=0$.

\begin{lemma}[The stationary action identity]\label{lem:actionidentity}
Put $r=1/t$, and let $x(r)$ be an algebraic critical branch of
$S=P/r-q$. A prime means total differentiation along the branch.
Then
\begin{equation}\label{eq:actionidentity}
 S'=-P(x(r))/r^2,\qquad q(x(r))=-S-rS'.
\end{equation}
If $x(r)$ is bounded, the polar part of $S$ is $\alpha/r$, where
$\alpha=P(x(0))$. If it escapes and \eqref{eq:nodenoframe} holds,
its polar part is neither zero nor $\alpha/r$ for any constant
$\alpha$.
\end{lemma}
\begin{proof}
Use a finite Puiseux cover so every branch is meromorphic. At a
critical point the chain-rule term $\nabla_x S\cdot x'$ vanishes,
giving both identities. For a monomial $c r^{-\kappa}$ of the
polar part of $S$, its contribution to $-S-rS'$ is
$(\kappa-1)c r^{-\kappa}$. Thus bounded $q(x(r))$ permits only
the pole $\alpha/r$ in $S$. Its coefficient is $P(x(0))$ when
$x$ is bounded.

For an escaping branch let $R=\|x\|$. The critical equation is
$\nabla P=r\nabla q$, so $\|\nabla P\|/R\longrightarrow0$.
The leading projective direction is a singular point of $F_3$.
The ordinary-node argument in the proof of the preceding coercive
lemma works at this one direction: write $x=sw+y$, $y=o(s)$,
in a fixed complementary plane. The injectivity there of
$H=\operatorname{Hess}F_3(w)$ forces $y$ bounded from
\[
 \nabla P/s=Hy+\nabla F_2(w)+O((1+\|y\|)^2/|s|)=o(1).
\]
Pairing with $w$ then gives $F_2(w)=0$. Condition
\eqref{eq:nodenoframe} gives $q_2(w)\ne0$, so $q(x(r))$ has a
pole: its leading quadratic term cannot be cancelled by its
lower-degree terms. Identity~\eqref{eq:actionidentity} now rules
out both claimed polar parts. This argument includes all escape
scales and all coincident actions.
\end{proof}

We specify the classical rank input before using it. This is a
convergent algebraic application of formal stationary phase, not
an extension to arbitrary nonconvergent coefficients or rings.

\begin{lemma}[Rank accounting for a tame algebraic pencil]
\label{lem:tamepencilrank}
For the present no-frame pencil, shrink a punctured $r$ disk to
avoid its finitely many exceptional parameters. Its top twisted
de Rham module has rank
\[
 N=\sum_{\text{critical points of }S_r}\mu.
\]
After a finite cover, the rank of its block with nonzero
exponential polar part $\gamma$ is the sum of these generic local
Milnor numbers over critical branches with action polar part
$\gamma$. Equal actions add their multiplicities. Consequently
the rank of the zero-exponential block is the sum of the generic
Milnor numbers of the branches with bounded action.
\end{lemma}
\begin{proof}
For fixed small nonzero $r$, the homogeneous quadratic term of
$S_r$ at a leading node $w$ is $F_2(w)/r-q_2(w)$. It is nonzero
at every node, by \eqref{eq:nodenoframe} and exclusion of finitely
many parameters. The ordinary-node coercivity lemma gives
$\|\nabla S_r\|\ge c_r\|x\|$ outside a compact set. The
critical locus is therefore finite: a positive-dimensional affine
algebraic component would be unbounded. Pham's tame Milnor
decomposition~\cite[Part~I, Sections~1--4]{Pham} gives the displayed
rank, retaining the Milnor number of every degenerate point.

Here are the checks needed to apply the direct-image rank formula.
Take $U=\C_r^*\times\C_x^3$, $f=r$, $g=S_r$, and the regular
algebraic module $\mathcal O_U$. Form the ordinary direct image
by $(f,g)$, followed by exponential integration in the $g$
coordinate. All these are algebraic maps; $g$ is regular on $U$.
The ordinary direct image is regular holonomic. Over the punctured
$r$ disk its relevant cohomology is the middle, reduced
Gauss--Manin module of the fibers of $S_r$. The other, constant
fiber cohomology has zero exponential integral in the affine $g$
line. Punctual parameter terms disappear on the punctured disk.

There are no extra finite-$g$ singular curves from spatial infinity.
The coercive estimate is locally uniform on a compact subset of
the punctured $r$ disk. For $g$ in a compact set, the identity
$P=r(g+q)$ on a fiber gives
\[
 \partial_r S_r=-(g+q)/r=O(1+\|x\|^2).
\]
Lifting $r$ variation by the normalized gradient therefore uses a
vector field of at most linear growth. Lifting $g$ variation uses
a field $O(1/\|x\|)$. These fields are complete outside a fixed
ball on each such parameter neighborhood. They trivialize the
outer fiber pairs, so all nonvertical finite-$g$ discriminant
curves come from affine critical values. Along such a curve the
conormal multiplicity of the middle Gauss--Manin module is the
sum of the local Milnor numbers: a transverse target disk computes
precisely the corresponding isolated vanishing cycles. This also
applies to persistent degenerate points and equal critical values.

First make a common finite Puiseux base change of all the
critical-value curves. Apply the nonramified case of Roucairol's
Theorem~0.1(2), proved in her Section~1.2, and Theorem~0.3
\cite{RoucairolFormal} to this identified ordinary direct image.
They give exactly the stated nonzero polar parts and ranks.
Count every graph after this base change, including distinct
conjugates with the same polar part. For example the two graphs
$g=r^{-1}\pm r^{1/2}$ must contribute twice after $r=\rho^2$;
retaining just their one polar expression would lose rank.
Their rank statement requires no separation of the constant terms
of equal polar parts. The additional separation hypothesis in her
Theorem~0.2 concerns its monodromy formula, which we do not use.
Vertical components at $r=0$ and the boundary $g=\infty$ do not
add polar branches in Theorem~0.1. Subtracting the nonzero-block
ranks from $N$ proves the last assertion. In particular we have
not inferred the regular rank from an Euler characteristic alone.
\end{proof}

\begin{proposition}\label{prop:ordinaryfinite}
Suppose the common ordinary-node pair satisfies
\eqref{eq:nodenoframe} and $P$ has finite affine critical set.
Then $P$ is not pure. More generally every original Gaussian
marked period is either zero or is not a finite sum of
polynomials in $t$ times exponentials $e^{\beta t}$.
\end{proposition}
\begin{proof}
Every bounded critical branch of $S_r$ limits to a critical point
of $P$. Disjoint small Milnor neighborhoods of these finitely many
points contain, for the deformation $P-rq$, total critical
multiplicity $\mu_p$ in the neighborhood of $p$. This is
conservation of the local intersection multiplicity of the
gradient, and does not require a Morse perturbation. By
Lemma~\ref{lem:actionidentity} its action has polar part
$P(p)/r$. Every escaping branch has nonzero polar part. Hence
Lemma~\ref{lem:tamepencilrank} identifies the entire zero-exponential
rank as
\begin{equation}\label{eq:ordinaryzerorank}
 \sum_{p:\,dP(p)=0,\,P(p)=0}\mu_p.
\end{equation}

We construct that many independent local formal functionals and
compute their monodromy over the unramified field $\C((r))$.
At an isolated critical point at least one of $u,v$ is nonzero.
Indeed if both vanish, criticality requires $h=e=f=0$, in which
case the whole $T$ axis is critical. In the chart $u\ne0$ put
\[
 s=uv+h,\quad R(u)=G(u,-h/u),\quad
 T'=T+\frac{G(u,(s-h)/u)-G(u,-h/u)}s,
\]
where $G=P-T(uv+h)$. The quotient is analytic, and exactly
$P=sT'+R(u)$. The coordinate Jacobian is $u\ne0$. At a critical
point $p$ the one-variable critical multiplicity of $R$ is
$\mu_p$. An analytic one-variable coordinate $y$, with
$w=T'$, therefore gives the exact identity
\[
 P=\alpha+sw+y^k,\qquad k=\mu_p+1,\quad\alpha=P(p).
\]
The volume Jacobian and $e^{-q}$ multiply every amplitude by an
analytic unit $V(y,s,w)$; no Gaussian mean or cross term has been
discarded.

For $j=1,\ldots,k-1$ define a formal local functional on an
analytic monomial $y^a s^b w^c$ as zero unless $b=c$ and
$a=j-1+\ell k$ for an integer $\ell\ge0$. Otherwise its value is
\begin{equation}\label{eq:ordinarycharacters}
 e^{\alpha/r}r^{1+j/k}
       (-1)^{b+\ell}b!\,(j/k)_\ell\,r^{b+\ell}.
\end{equation}
Evaluate a polynomial form after multiplying its transformed
amplitude by $V$. At each power of $r$ only finitely many Taylor
terms contribute. The recurrences
\[
 I_{a+k}=-\frac{a+1}{k}r I_a,\qquad
 J_{b+1,b+1}=-(b+1)rJ_{b,b}
\]
prove spatial integration by parts and parameter horizontality.
Thus these are functionals on the full twisted quotient, not on
an assumed cyclic image. They exist even if the corresponding
asymptotic series diverge; their input coefficients here are the
Taylor series of the stated analytic coordinates.

They are independent. Prescribe, at each $p$, the jets
$1,y,\ldots,y^{k-2}$ on the line $s=w=0$ and zero jets at the
other points, using polynomial representatives of finite local
jets and the Chinese remainder theorem. After dividing column
$j$ by $e^{\alpha/r}r^{1+j/k}$, the leading evaluation matrix is
the jet map multiplied by the analytic unit $V$. It is triangular
with nonzero diagonal. Its higher terms are in $r\C[[r]]$.
For $\alpha=0$ these independent characters fill exactly the rank
in \eqref{eq:ordinaryzerorank}; there is no unaccounted regular
piece. Their formal monodromy eigenvalues, or their inverses if
one uses the dual module, are
\[
 \exp(2\pi i j/k),\qquad 1\le j<k.
\]
None is one. This conclusion is over $\C((r))$. We do not replace
a noninteger residue by its integer multiple after ramification.
The unique zero-exponential projector commutes with descent, and
the displayed unramified coefficient matrices identify its whole
block. Every other block has a nonzero exponential polar part.
The ambient module therefore has no trivial composition factor.

For completeness this excludes the formal Gaussian unit without
choosing an actual analytic period. Formal Gaussian integration
annihilates $\partial_i+tP_i-q_i$ coefficientwise. If its value
on a fixed marker $A$ were a nonzero constant, restriction to the
rational differential submodule generated by $A$ would give a
surjection to the trivial module: apply the functional to every
rational relation between $A,\nabla_tA,\ldots$. A relation has
zero constant coefficient because all positive derivatives of
that constant vanish. Exact formal extension at infinity would
preserve this composition factor, contradicting what was proved.

Finally a nonzero finite exponential-polynomial can be reduced
to $c e^{\beta t}$ by a constant-coefficient differential
operator. On an original marker this operator replaces $A$ by
the fixed polynomial marker $T(P)A$. Replace $P$ by $P-\beta$
and divide the marker by $c$. The hypotheses and finite critical
set are unchanged, so the same contradiction applies. Taking
$A=1$ proves the purity exclusion.
\end{proof}

\subsection{An exhaustive critical-set split}

The finite-critical hypothesis in Proposition~\ref{prop:ordinaryfinite}
has a short exhaustive coefficient test. This prevents a generic
argument from silently discarding the positive-dimensional cases.

If $h\ne0$, the critical locus lies on $uv=-h$, and the equations
above identify it with $R'(u)=0$ in $\C^*$. The Laurent polynomial
\[
 R(u)=a u^3+b u^2+e u+k-fh/u+ch^2/u^2-dh^3/u^3
\]
is nonconstant unless $a=b=c=d=e=f=0$. A nonzero Laurent
derivative has only finitely many roots in $\C^*$. In the constant
case the critical set is the smooth conic $uv=-h$, $T=0$.

If $h=0$, outside the center the critical equations are
\[
 \begin{array}{lll}
 v=0,\ u\ne0:&3au^2+2bu+e=0,&T=-f/u,\\
 u=0,\ v\ne0:&3dv^2+2cv+f=0,&T=-e/v.
 \end{array}
\]
The center $u=v=0$ is critical exactly when $e=f=0$, and then
the whole $T$ axis is critical. Consequently the following table
is exhaustive and disjoint:
\begin{center}\small
\begin{tabular}{@{}p{0.51\textwidth}p{0.39\textwidth}@{}}
\toprule
Condition & Conclusion\\
\midrule
$h\ne0$, $(a,b,c,d,e,f)\ne0$ & Finite critical set; excluded above\\
$h\ne0$, $a=b=c=d=e=f=0$ & Critical conic; Corollary~\ref{cor:ordinaryconics}\\
$h=0$, $e=f=0$ & Center axes; Proposition~\ref{prop:axisexclusion}\\
$h=0$, $(e,f)\ne0$, $a=b=e=0$ & Critical hyperbola; Corollary~\ref{cor:ordinaryconics}\\
$h=0$, $(e,f)\ne0$, $d=c=f=0$ & Symmetric hyperbola; same corollary\\
$h=0$, $(e,f)\ne0$, $(a,b,e)\ne0$, $(d,c,f)\ne0$
 & Finite critical set; excluded above\\
\bottomrule
\end{tabular}
\end{center}
The two hyperbola rows cannot overlap because $(e,f)\ne0$.
All isolated degeneracies in the finite rows have already been
included through their Milnor multiplicities. The conic and hyperbola rows are proved in the next subsection.
The center-axis row is proved in Proposition~\ref{prop:axisexclusion},
using a full formal frame or the uniform actual-cut criterion.

\subsection{A critical conic in every polynomial degree}
\label{sec:hyperbolicsuspension}

The conic and hyperbola rows of the preceding table are manifestations
of the same calculation. It also works when one of the cubic terms is
replaced by a polynomial of arbitrary degree. We prove that stronger
statement to make clear what the mechanism actually uses.

\begin{theorem}\label{thm:hyperbolicsuspension}
Retain a full nondegenerate affine Gaussian quadratic
\[
 q=\tfrac12(Ax^2+By^2+mw^2+2Cxy+2Dxw+2Eyw)
                     +\ell_x x+\ell_y y+\ell_w w+q_0.
\]
Let $AB\ne0$ and
\[
 P=wxy+F(w)+k,\qquad F(0)=0,\quad F'(0)=h\ne0.
\]
Allow any polynomial degree $d_F=\deg F\ge1$, requiring $m\ne0$
when $d_F=1$. Then $P$ is not pure. Every original marked Gaussian
period is either zero or is not a finite exponential-polynomial.
The generic top twisted rank is $d_F+3$ for $d_F\ge2$ and five for
$d_F=1$. No Morse assumption on the Laurent phase obtained by
restricting $q$ to the conic, or on the isolated critical points of
$P$, is needed.
\end{theorem}

The nonzero quadratic poles at both ends of the conic are essential
to the rank-four construction below. The condition $AB\ne0$ states
this explicitly. Nondegeneracy of the full Gaussian alone would not
imply it. This theorem concerns complex polynomials and valued-field
formal frames; it does not assert an all-ring extension.

\begin{proof}
We first justify the full rank. Fix a generic $t\ne0$ and set
$S=tP-q$. Its gradient has a linear lower bound at complex spatial
infinity. To see this, suppose along an unbounded sequence
$\|\nabla S\|=o(R)$, where $R=\|(x,y,w)\|$. If $w$ stays
bounded, pass to a limit $w_0$ and normalize $(x,y)$ by $R$.
The first two gradient equations give
\[
 -AX+(tw_0-C)Y=0,\qquad (tw_0-C)X-BY=0.
\]
At least one of $X,Y$ is nonzero, and $AB\ne0$ makes both
nonzero. But the last gradient equation, divided by $R^2$, gives
$tXY=0$, a contradiction.

If $w$ is unbounded, the first two equations have a matrix with
inverse $O(1/|w|)$. Hence
$x,y=O(1)+o(R)/|w|$. First $R\asymp |w|$, then $x,y=O(1)$.
The last equation is therefore
\[
 tF'(w)-mw+O(1)=o(|w|).
\]
This is impossible for $d_F\ge3$, for $d_F=2$ with
$2t[w^2]F-m\ne0$, and for $d_F=1$ with $m\ne0$.
The same proof is uniform on compact subsets of these coefficient
conditions. Thus all phase-critical points stay in one ball during
any compact coefficient deformation inside that domain. They are
isolated, and their total Milnor number is constant by conservation
of local gradient-intersection multiplicity. Pham's tame Milnor
decomposition~\cite[Part~I, Sections~1--4]{Pham} identifies that
number with the full twisted rank.

For fixed $t$, the coefficient domain just used is the complement
of finitely many complex algebraic hypersurfaces, and is path
connected. We may compute the number at a generic diagonal
Gaussian with zero means, retaining $ABm\ne0$, and a generic $F$.
The critical equations there are
\[
 twy=Ax,\qquad twx=By,\qquad t(xy+F'(w))=mw.
\]
There are $d_F-1$ simple solutions with $x=y=0$ if $d_F\ge2$,
and one if $d_F=1$. The other solutions have
$w=\pm\sqrt{AB}/t$. At each of these two values, $x,y$ are
proportional and their nonzero product is prescribed by the last
equation; generically this gives two simple solutions. The counts
are $d_F+3$ and five, respectively. Uniform coercivity and the
connected deformation justify this specialization; a convenient
critical-point count by itself would not justify the rank for a
degenerate member.

We now construct a full formal frame. Set $r=t^{-1/2}$ and write
$F(w)=hw+bw^2+\sum_{j\ge3}f_jw^j$. Near the entire critical
conic $w=0$, $xy=-h$, use the exact coordinates
\[
 y=(rV-h)/x,\qquad w=rU.
\]
They give
\begin{align}
 tP-q&=k/r^2+UV+bU^2+\psi(x)+E_r(x,V,U),
                  &E_r\in r\C[x,x^{-1},V,U][r],\label{eq:conicphase}\\
 dx\,dy\,dw&=r^2\,\frac{dx}{x}\,dV\,dU,\label{eq:conicvolume}\\
 \psi(x)&=-\frac A2x^2-\frac{Bh^2}{2x^2}
                -\ell_x x+\frac{h\ell_y}{x}+Ch-q_0.\notag
\end{align}
All Gaussian cross terms and means are in $\psi$ or $E_r$.
The normal quadratic $UV+bU^2$ has Hessian determinant $-1$ for
every $b$. Its normalized formal Gaussian functional is well
defined and takes one to one.

The logarithmic Laurent quotient for $\psi$ has rank four, with
basis $1,x,x^2,x^3$. Here is a direct proof, including independence.
Its relation on $x^j$ is
\[
 jx^j-Ax^{j+2}-\ell_xx^{j+1}
             -h\ell_yx^{j-1}+Bh^2x^{j-2}=0.
\]
The top coefficient eliminates all exponents at least four without
creating a negative exponent; the bottom eliminates all negative
exponents without creating an exponent above three. If the image
of a nonzero Laurent polynomial under $x\partial_x+x\psi_x$
were supported in $[0,3]$, its highest exponent would be at most
one and its lowest at least two, which is impossible.

Let $\lambda_0,\ldots,\lambda_3$ be the dual basis on this
quotient, tensored with the normal Gaussian functional. For a
polynomial amplitude $a(x,y,w)$, define
\begin{equation}\label{eq:coniccharacters}
 L_i(a)=e^{k/r^2}r^2\lambda_i\!\left(
    e^{E_r}a(x,(rV-h)/x,rU)\right).
\end{equation}
Every coefficient is a finite Laurent-polynomial calculation, so
this is a well-defined formal series. Multiplication by $e^{E_r}$
conjugates the fixed twisted differential to the perturbed one.
Consequently these four functionals annihilate all spatial exact
forms. The moving volume is retained in the factor $r^2$.
For an explicit parameter check the coordinate velocity in the
original variables is
\[
 Z=\left(0,\frac{xy+h}{rx},\frac wr\right),
                   \qquad \operatorname{div}Z=2/r.
\]
Cartan's identity $\mathcal L_Z=d\iota_Z+\iota_Zd$ shows that
the pullback derivative differs from the original parameter
connection only by a spatial exact form and this volume
derivative. Thus \eqref{eq:coniccharacters} are horizontal for
the original full module. In particular, for every fixed
polynomial $a$,
\begin{equation}\label{eq:conicseedorder}
 L_i(a)\in e^{kt}r^2\C[[r]].
\end{equation}
This bound holds for all original markers, not just for the unit.

Suppose first $d_F\ge2$. The other critical points of $P$ are
$x=y=0$, $F'(w_p)=0$. Each $w_p$ is nonzero because $h\ne0$.
If it has multiplicity $\mu_p$, analytic coordinates give
\[
 P=k+F(w_p)+s z+y_0^{\mu_p+1}.
\]
Indeed absorb the nonvanishing factor $w$ into $x$ or $y$, then
use the one-variable power coordinate for $F$ at $w_p$.
The explicit characters \eqref{eq:ordinarycharacters}, including
their integration-by-parts and time recurrences, apply to this
exact local model. Their prefactors are
\[
 e^{(k+F(w_p))t}t^{-1-j/(\mu_p+1)},
                    \qquad 1\le j\le\mu_p.
\]
The Gaussian and coordinate Jacobian are analytic units in that
construction. These characters retain multiple roots and equal
critical values. Their total number is $\sum_p\mu_p=d_F-1$.

The $d_F+3$ characters just written are independent on polynomial
amplitudes. Put $B_0(w)=\prod_p(w-w_p)^{\mu_p}$. For the four
conic rows use $x^jB_0(w)$, $0\le j\le3$. Their leading
conic matrix is $B_0(0)I_4$. Their leading finite jets vanish.
For the remaining rows choose polynomials in $w$ with value zero
at $w=0$ and the separate Hermite jets of orders
$0,\ldots,\mu_p-1$ at the nonzero $w_p$. The Chinese remainder
theorem supplies them. Their conic values have an additional
positive power of $r$. Their normalized finite evaluation matrix
is the invertible local jet map, multiplied by analytic units.
After dividing each character column by its scalar prefactor,
the leading matrix is therefore block diagonal and invertible.
Since its size is the proved full rank, these are a full formal
frame. No uncomputed image kernel or infinity block remains.

For $d_F=1$ there is instead one further character. In the exact
coordinates
\[
 x=X/t,\qquad y=Y/t,\qquad w=(h/m)t+W,
\]
the phase is
\begin{align*}
 tP-q={}&\frac{h^2}{2m}t^2+(k-h\ell_w/m)t+H(X,Y,W)+O(t^{-1}),\\
 H={}&\frac hm XY-\frac hm(DX+EY)-\frac m2 W^2-\ell_w W-q_0.
\end{align*}
The correction is polynomial in $t^{-1},X,Y,W$, with positive
$t^{-1}$ order. The Hessian determinant of $H$ is $h^2/m\ne0$.
Formal Gaussian evaluation gives a horizontal character with
nonzero quadratic action and volume prefactor $t^{-2}$.
To check independence set $w_*=(h/m)t$. Use
$x^j(1-w/w_*)$, $0\le j\le3$, and $w/w_*$ as five polynomial
amplitudes with rational parameter coefficients. The former
have the original conic leading matrix and zero leading escape
values; the latter has zero leading conic values and nonzero
leading escape value. This gives the full rank-five frame.

In either degree case, after finite ramification the frame splits
the module into rank-one factors. Every zero-exponential
functional has positive order on every original polynomial
marker, by \eqref{eq:conicseedorder} and the displayed finite
prefactors. Lemma~\ref{lem:positiveformal} rules out a nonzero
constant original marked period. A constant-coefficient operator
reduces any nonzero finite exponential-polynomial to $c e^{\beta t}$;
it replaces the marker by another original polynomial in $P$.
Shifting $k$ by $-\beta$ and normalizing the marker gives the same
contradiction. This proves the theorem.
\end{proof}

\begin{corollary}\label{cor:ordinaryconics}
All the critical-conic and critical-hyperbola rows in the ordinary-node
table are not pure. Only its center-axis row remains open here.
\end{corollary}
\begin{proof}
For $h\ne0$ and $a=b=c=d=e=f=0$ the residual is
$T(uv+h)+k$. Apply the $d_F=1$ case with $w=T$, $x=u$, $y=v$.
The no-frame conditions give $ABm\ne0$.

For $h=0$, $a=b=e=0$, $f\ne0$, permute the coordinates and
write the residual as
\[
 v(Tu+d v^2+c v+f)+k.
\]
Take $w=v$, $x=T$, $y=u$, and
$F(w)=d w^3+c w^2+f w$. The quadratic coefficients along $T,u$
are nonzero: the original $T$ square coefficient is nonzero, and the $u$
node has $F_2=0$, so its Gaussian square coefficient is nonzero
by the no-frame condition. If this $F$ is linear, the remaining
$v$ node also has $F_2=0$, giving the extra $w$ square coefficient
required in degree one. The symmetric hyperbola row is identical
after swapping $u,v$. These are all the indicated rows.
\end{proof}

\subsection{A uniform cut criterion in any degree}
\label{sec:uniformcutdecay}

Intersecting critical curves need not admit a simple formal normal
form. The following criterion avoids choosing one. Its hypothesis
concerns the full phase $tP-q$, uniformly in large complex $t$;
bounded critical points alone would not suffice.

\begin{theorem}\label{thm:uniformcuts}
Let $P,q\in\C[x_1,\ldots,x_n]$, with $P$ nonconstant. Suppose
there are fixed $c,R_0,t_0>0$ such that
\begin{equation}\label{eq:uniformphasegradient}
 \|\nabla_x(tP-q)(x)\|\ge c\|x\|
       \quad(|t|\ge t_0,\ \|x\|\ge R_0).
\end{equation}
Along every fixed ray $t=Re^{i\theta}$, the full rapid-decay
homology has a flat cut basis $\Gamma_i(t)$, attached to the
critical values $c_i(t)$ of $tP-q$, such that for every fixed
polynomial $a$,
\begin{equation}\label{eq:normalizedcutdecay}
 e^{-c_i(t)}\int_{\Gamma_i(t)}a(x)e^{tP-q}\,
                 dx_1\cdots dx_n\longrightarrow0.
\end{equation}
Each $c_i(t)=\alpha_i t+O(1)$, where $\alpha_i$ is a critical
value of $P$. If all critical values of $P$ lie in a closed
half-plane through zero, a supporting $t$ ray makes every flat
period of every such fixed form tend to zero.
\end{theorem}

We give the completeness, tail and limiting-form arguments separately.
None follows just from counting the phase-critical points.

\begin{proof}
For every sufficiently large $t$, \eqref{eq:uniformphasegradient}
makes the phase-critical set finite and contained in one fixed
ball. It also gives complex tameness at infinity. Pham's tame
Milnor decomposition~\cite[Part~I and Appendix]{Pham} identifies
the complete middle rapid-decay homology with the local vanishing
homologies carried along cuts from its critical values. The rank
is the sum of their Milnor numbers, including degenerate points.
If this rank is zero the conclusion is immediate. Otherwise
retain all these local generators; identical critical values
are grouped, not discarded.

On the chosen ray, the critical points and values are algebraic
functions of $R$, after a finite cover. Enlarge $R_1$ so branches
and multiplicities stabilize. Choose $|\eta|=1$ with
$\Re\eta\ge1/2$ so the cuts
\[
 S_t(x)=c_i(t)-\eta s,\qquad s\ge0,\quad S_t=tP-q,
\]
are disjoint for $R>R_1$, except at identical grouped endpoints.
There are only finitely many forbidden choices: every nonzero
difference of two algebraic critical values has one limiting
direction. The cut angle is independent of $\theta$. In
particular a supporting $t$ ray need not be a generic ray.

Along a cut, normalized complex gradient transport is
\[
 x_s=-\eta\,\overline{\nabla S_t}/\|\nabla S_t\|^2,
               \qquad S_t(x(s))_s=-\eta.
\]
Outside the fixed ball,
$|d\|x\|^2/ds|\le2/c$. All local vanishing generators start
in that ball. With fixed sufficiently large $A$ and $B>2/c$,
they therefore have representatives inside the carrier
\begin{equation}\label{eq:uniformcarrier}
 \|x\|^2\le A+Bs.
\end{equation}
At its outer boundary the transport points strictly inward
relative to this inequality. No trajectory escapes at finite
cut length, and the cuts avoid all other critical values.

We must control a single family of representatives. The gradient
flow itself need not be semialgebraic. Instead compactify $x$
by $x/\sqrt{1+\|x\|^2}$ and $s$ by $s/(1+s)$. Take the closure
of each constrained cut carrier and mark its rapid boundary
$G$, where $s=\infty$. Inequality~\eqref{eq:uniformcarrier}
forbids spatial infinity at finite $s$. For fixed $R$, the
relative chains of this compact pair are precisely locally
finite cut chains with their ends allowed only at rapid
infinity. This follows also by a compatible triangulation of
the compact pair: its relative simplices restrict to the usual
locally finite chain model in the open carrier. The exponential
there is exactly $e^{c_i-\eta s}$.

For the full unconstrained cut carrier, mark and remove any
spatial boundary at finite $s$; it is not a rapid boundary.
The constrained carrier, the full carrier, their rapid boundaries
and this excluded boundary form a finite family of semialgebraic
subsets in the same compactification. Apply the compatible
finite-family form of Hardt triviality to projection onto $R$.
It gives one trivialization on an unbounded interval, preserving
these subsets and their inclusions. The precise form is stated
in Drusvyatskiy--Lewis~\cite[Theorem~2.25]{DrusvyatskiyLewis};
it is Hardt's theorem. Semialgebraic triangulation and the
chain-current boundary rule are as in
Hardt--Lambrechts--Turchin--Voli\'c~\cite[Sections~2.1 and~3,
especially Theorem~3.5]{HLTV}.

Gradient transport already showed that the constrained inclusion
contains every local vanishing generator. At one parameter choose
finitely many integral relative chains whose images are a basis
over $\C$ of the full rapid-decay space. A nonprimitive integral
image causes no problem over $\C$. Transport these chains by
the compatible trivialization. They form a jointly semialgebraic
family with fixed integer multiplicities and supports satisfying
\eqref{eq:uniformcarrier}. Their images are a flat basis of the
actual homology local system: on each compact parameter interval
the transport preserves the relative pairs and is proper on
every finite phase band. Thus this is a completeness statement
for actual cycles, not an identification of presumed formal blocks.

Here is the needed uniform mass bound. In a fixed semialgebraic
family of $d$-chains with bounded multiplicities, the mass inside
a ball of radius $L\ge1$ is at most $C L^d$, uniformly in the
parameters. Partition tangent planes by their largest coordinate
$d$-minor. The Cauchy--Binet identity gives a lower bound
$\binom Nd^{-1/2}$ for some such minor in ambient dimension $N$.
The area formula then bounds that piece by its coordinate
projection volume times the number of preimages. Semialgebraic
uniform finiteness bounds those finite fibers independently of
the parameters; see, for example,
\cite[Theorem~2.32]{DrusvyatskiyLewis}. Sum the finitely many
pieces and multiplicities. Apply the same argument to boundary
chains. This proves uniformity, rather than merely pointwise
finite volume.

Lift the cut chains to $(x,s)$. Their bands $j\le s\le j+1$
lie in balls of radius $O(1+j)$ by
\eqref{eq:uniformcarrier}. Their fixed semialgebraic descriptions
and the mass estimate give, with constants independent of $R$,
\begin{equation}\label{eq:uniformcuttail}
 \left|e^{-c_i(t)}\int_{s\ge L}a e^{S_t}\,dx_1\cdots dx_n\right|
 \le C\sum_{j\ge\lfloor L\rfloor}(1+j)^M e^{-j/2}.
\end{equation}
The polynomial amplitude changes only $M$. This tail tends to
zero uniformly in $R$.

It remains to bound a fixed phase band. A bounded critical branch
$x(r)$, $r=1/t$, satisfies
\[
 q(x(r))=-c_i(r)-r c_i'(r).
\]
At a stationary point the chain rule gives $c_i'(r)=-P(x(r))/r^2$;
eliminating $P$ gives this identity. Since $q(x(r))$ is
bounded, its Puiseux expansion forces the only possible polar
term of $c_i$ to be $\alpha_i/r$. Its coefficient is
$\alpha_i=P(x(0))$, and $\nabla P(x(0))=0$. Hence
$c_i(t)=\alpha_i t+O(1)$ even at critical collisions.

Choose a generic fixed truncation level $L$. The truncated lifted
chains have compact supports and uniformly bounded mass and
boundary mass. The latter follows from the same semialgebraic
estimate applied to their slices; after a finite stratification
there is an arbitrarily large sequence of such levels. By the
Federer--Fleming compactness theorem, every sequence $R\to\infty$
has a subsequence converging to an integral $n$-current. We use
the version with both mass and boundary-mass bounds,
Simon~\cite[Chapter~6, Section~3, Theorem~3.15]{SimonGMT}.
Its rectifiability conclusion is essential.

On the lifted chain,
\[
 P(x)-\alpha_i=
       \frac{c_i(t)-\alpha_i t+q(x)-\eta s}{t}.
\]
Every limiting current is therefore supported in
$\{P=\alpha_i\}\times[0,L]$. Stratify this complex algebraic
level set into smooth complex strata. Almost everywhere for the
rectifiable limiting current, its tangent plane projects into
the tangent of one of these strata. Each has complex dimension
less than $n$. The holomorphic $n$-form $dx_1\cdots dx_n$ thus
vanishes on that tangent plane. Lower-dimensional strata carry
no $n$-mass; the same argument covers singular strata of larger
real dimension. Testing with the fixed smooth form
$e^{-\eta s}a(x)dx_1\cdots dx_n$, with a compact cutoff, gives
zero on every subsequential limit. Consequently the normalized
integral on $0\le s\le L$ tends to zero. This does not assert
that the real volume of the chains tends to zero.

First let $R\to\infty$, then let $L\to\infty$ through the
allowed levels in \eqref{eq:uniformcuttail}. This proves
\eqref{eq:normalizedcutdecay}. Finally choose $\theta$ with
$\Re(e^{i\theta}\alpha)\le0$ for all critical values $\alpha$
of $P$. The factors $e^{c_i(t)}$ are then bounded, so every basis
period tends to zero. A flat cycle is a constant linear
combination of this full basis, proving the last assertion.
\end{proof}

\begin{corollary}\label{cor:twocriticalvalues}
Let $q$ be a full affine Gaussian quadratic and assume
\eqref{eq:uniformphasegradient}. If $P$ has at most two critical
values, every original marked Gaussian period is either zero or
is not a finite exponential-polynomial. In particular $P$ is
neither pure nor safe.
\end{corollary}
\begin{proof}
Any two complex numbers lie in a closed half-plane through zero.
For two nonzero values normalize them to $z_1,z_2$ on the unit
circle. If $z_1+z_2\ne0$, use the direction
$-\overline{z_1+z_2}/|z_1+z_2|$; if their sum is zero, use
$i\overline{z_1}$. Zero or missing values add no condition.

A nonzero constant original marked period would, by the rational
seed restriction and actual-period extension in
Lemma~\ref{lem:seedfunctional}, be represented by a flat period
of that fixed polynomial form. Its generic exponential system
has $f=t$, $g=tP-q$ on smooth affine space and regular input
$\mathcal O$; $f,g$ are algebraically independent because $P$
is nonconstant. Remove the finite exceptional parameter set of
the Hien--Roucairol comparison. Theorem~\ref{thm:uniformcuts}
makes the resulting constant period tend to zero, a contradiction.
No convergence of the formal moment series was assumed.

A constant-coefficient operator reduces a nonzero finite
exponential-polynomial to $c e^{\beta t}$ while replacing its
original marker by another fixed polynomial in $P$ times that
marker. Shift $P$ by $-\beta$ and normalize. Its gradient is
unchanged and its critical values are still at most two. The
same contradiction applies. Since $\Phi_1(0)=1$, it cannot be
a polynomial, proving also the safety exclusion.
\end{proof}

\subsection{Completing the ordinary-node center axes}
\label{sec:ordinaryaxes}

The only remaining ordinary-node row is
\begin{equation}\label{eq:ordinaryaxes}
 P=Tuv+a u^3+d v^3+b u^2+c v^2+k.
\end{equation}
Keep the full $q$ of Lemma~\ref{lem:ordinarynormal}, with
$m\ne0$. Under the no-frame assumption the other nodes satisfy
\[
 a=b=0\ \Longrightarrow\ A\ne0,\qquad
 d=c=0\ \Longrightarrow\ B\ne0.
\]
We use a full formal frame when $bc\ne0$ and the preceding
uniform cut criterion when $bc=0$. These branches are disjoint
and include every zero face of the four coefficients.

\begin{lemma}\label{lem:axisuniformgradient}
For \eqref{eq:ordinaryaxes} with these Gaussian conditions,
\eqref{eq:uniformphasegradient} holds.
\end{lemma}
\begin{proof}
Otherwise take $\|x\|\to\infty$, $|t|\ge t_0$ and
$\nabla(tP-q)=o(\|x\|)$. Division by $t\|x\|^2$ shows that
a limiting projective direction is a leading node: the $T$ axis,
also the $u$ axis if $a=0$, and the $v$ axis if $d=0$.

Near any such node use height $z$ and transverse coordinates
$x,y=o(z)$. The phase polynomial has the form
\[
 P=zxy+\beta z^2+H(x,y)+k,
\]
where $H$ has only terms of degrees two and three. There are
no transverse linear terms. The transverse equations, divided
by $tz$, give
\[
 |x|+|y|\le C/|t|+o(|x|+|y|)+o(1/|t|).
\]
The terms of $\nabla H/z$ are absorbed because $(x,y)=o(z)$;
the Gaussian cross terms divided by $tz$ are bounded by the
same quantities. Thus $x,y=O(1/|t|)$. The height equation is
\[
 \partial_z(tP-q)=(2\beta t-M)z+O(1),
\]
where $M$ is the Gaussian square coefficient along this node.
If $\beta\ne0$, $|2\beta t-M|$ is uniformly bounded below
for all sufficiently large $|t|$. If $\beta=0$, the no-frame
condition gives $M\ne0$. This contradicts the assumed small
gradient ratio. The proof also covers bounded-$t$ subsequences
within the chosen large-$|t|$ domain and therefore proves the
required uniform estimate.
\end{proof}

\begin{lemma}\label{lem:axisrank}
The full generic twisted rank of \eqref{eq:ordinaryaxes} is
\[
 N=5+\mathbf1_{a\ne0}+\mathbf1_{d\ne0}.
\]
\end{lemma}
\begin{proof}
The three homogenized quadratic gradient equations have total
intersection length eight in $\mathbb P^3$. Their points at
infinity are exactly the leading ordinary nodes, in number
$3-\mathbf1_{a\ne0}-\mathbf1_{d\ne0}$. Each has length one.
Indeed its local Jacobian has the two transverse columns of
$\operatorname{Hess}F_3(w)$ and the column $\nabla(tF_2-q_2)(w)$.
The first two have rank two and are annihilated by pairing with
$w$; the last has nonzero pairing
$2(tF_2-q_2)(w)$. This value is nonzero for large generic $t$
by the stated conditions. Thus these infinity intersections are
simple. The uniform gradient estimate gives a finite affine
critical set; there are no other components. B\'ezout and
Pham's tame comparison give the asserted full rank.
\end{proof}

\begin{proposition}\label{prop:axisexclusion}
Every phase \eqref{eq:ordinaryaxes} under the stated conditions
is not pure. Every original marked Gaussian period is zero or
is not a finite exponential-polynomial.
\end{proposition}
\begin{proof}
First let $bc\ne0$. Set $r=t^{-1/2}$, $u=rU$, $v=rV$, leaving
$T$ unchanged. Exactly,
\begin{align*}
 tP-q={}&k/r^2+F_0(U,V,T)+E_r(U,V,T),\\
 F_0={}&TUV+bU^2+cV^2-mT^2/2-\ell_TT-q_0,\\
 E_r={}&r(aU^3+dV^3-DUT-EVT-\ell_uU-\ell_vV)\\
       &-r^2(AU^2+BV^2+2CUV)/2,\\
 du\,dv\,dT={}&r^2dU\,dV\,dT.
\end{align*}
The full polynomial twisted quotient for $F_0$ has rank five,
including its degenerate parameter values. Here is a separate rank
proof. Suppose $\nabla F_0=o(\|(U,V,T)\|)$ on an unbounded
sequence. If $T$ stays bounded, the first two equations force
both normalized $U,V$ to be nonzero; the last, divided by the
squared norm, gives their product zero, a contradiction. If
$T$ is unbounded, the first two equations have inverse
$O(1/|T|)$, giving $U,V=o(1)$ and $\|(U,V,T)\|\asymp|T|$.
The last equation then has leading term $-mT$, again a
contradiction. This coercivity is uniform on compact sets with
$bcm\ne0$. The coefficient domain is connected, so conservation
of critical intersection multiplicity and Pham's tame comparison
make its rank constant. At generic $\ell_T$, its critical
equations give one point with $U=V=0$ and four with
$T^2=4bc$, all simple. Hence the rank is five, also when points
merge at the exceptional values of $\ell_T$.

Choose any polynomial basis $B_1,\ldots,B_5$ of this fixed
quotient and its dual $\lambda_1,\ldots,\lambda_5$. Then
\[
 L_j(a_0)=e^{kt}r^2\lambda_j\bigl(e^{E_r}a_0(rU,rV,T)\bigr)
\]
defines five formal horizontal functionals on the original full
module. Coefficients are finite polynomial evaluations at every
order. Spatial exactness is the fixed-phase conjugacy; time
horizontality retains the displayed $r^2$ volume and the
coordinate velocity $(u/r,v/r,0)$, whose divergence is $2/r$.
Cartan's formula makes the remaining velocity term spatially exact.
For every original fixed polynomial marker their values lie in
$e^{kt}r^2\C[[r]]$.

The other $P$-critical points are the off-axis points
$p_u=(-2b/(3a),0,0)$ if $a\ne0$ and
$p_v=(0,-2c/(3d),0)$ if $d\ne0$. They are Morse, with Hessian
determinants $8b^3/(9a^2)$ and $8c^3/(9d^2)$. Holomorphic Morse coordinates reduce their local phases to a
nondegenerate quadratic. Expanding the analytic Gaussian amplitude
and volume unit and applying the quadratic Wick recurrence gives
all-order horizontal characters. These have prefactors $e^{P(p)t}r^3$, multiplied by power series in
$r^2$ on fixed markers. Their values may coincide with each
other or with $k$; no separation hypothesis is used.

Here is an explicit completeness test. Let $M$ bound the degree
of the chosen $B_j$ in $U,V$. Substitute $U=u/r,V=v/r$ and
multiply each $B_j$ by
\[
 C(u,v)=\prod_{p_u}(1-u/u_p)^K
                      \prod_{p_v}(1-v/v_p)^K,\qquad K>M.
\]
An empty product is one. At the central fixed model this is
$1+O(r)$, so the normalized five-by-five matrix is the identity
at leading order. At each off-axis point its zero of order $K$
overcomes the worst $r^{-M}$ growth of the substituted basis;
the normalized cross values tend to zero. For the off-axis rows
use $(u/u_p)^K$ or $(v/v_p)^K$. They have zero central leading
values, unit value at their own Morse point and zero leading
value at the other. The resulting full matrix has invertible
block diagonal leading term. Its size is exactly the rank of
Lemma~\ref{lem:axisrank}. These are therefore a full formal
frame, including all fixed-model degeneracies and equal actions.

Every zero-exponential original-marker seed has positive order,
either $r^2$ or $r^3$. Lemma~\ref{lem:positiveformal} excludes
a nonzero constant original marked period. Differential
isolation and shifting $k$ then exclude every nonzero finite
exponential-polynomial as before.

It remains to treat $bc=0$. The critical equations of $P$ show
that its values are among
\[
 k,\qquad k+4b^3/(27a^2)\quad(ab\ne0),\qquad
 k+4c^3/(27d^2)\quad(dc\ne0).
\]
At a missing isolated point omit its displayed value. Any extra
critical axes have value $k$. Since $bc=0$, there are at most
two values, including every intersecting-axis or degenerate
case. Lemma~\ref{lem:axisuniformgradient} and
Corollary~\ref{cor:twocriticalvalues} apply. This completes both
disjoint branches and proves the proposition.
\end{proof}

\begin{corollary}\label{cor:ordinarycomplete}
Frame extraction is proved for all three ordinary-node leading
types: irreducible nodal cubics, transverse line--conic unions,
and three nonconcurrent lines. In fact none of these leading
types permits a pure polynomial.
\end{corollary}
\begin{proof}
If all leading nodes have $F_2\ne0$, use the coercive-node
criterion. Otherwise an isotropic zero-$F_2$ node supplies the
frame directly. In the absence of such a frame, the common
normal form applies and satisfies the no-frame conditions at
every other node. The exhaustive critical-set table is now
entirely covered by Proposition~\ref{prop:ordinaryfinite},
Corollary~\ref{cor:ordinaryconics}, and
Proposition~\ref{prop:axisexclusion}. Thus a pure polynomial
must have a frame.

In that frame, the quadratic coefficient of $W$ in the cubic
leading form is the tangent quadratic at the chosen ordinary
node. It has rank two, so its discriminant is nonzero.
Proposition~\ref{prop:nondeglog} excludes purity in that branch
as well. This proves the stronger assertion.
\end{proof}


\subsection{The geometric case split for frame extraction}

For a cubic of degree exactly three, either the frame test
\eqref{eq:frame} succeeds, or it does not. In the latter case,
the projective curve $F_3=0$ belongs to exactly one of these types:
\begin{center}\small
\begin{tabular}{@{}p{0.33\textwidth}p{0.56\textwidth}@{}}
\toprule
Leading curve & Proof or remaining no-frame exclusion\\
\midrule
Smooth cubic & Proved by Corollary~\ref{cor:smoothleading}; no pure phase\\
Irreducible nodal cubic & Completely excluded by Corollary~\ref{cor:ordinarycomplete}\\
Irreducible cuspidal cubic & Smooth parabola, separated lines and double-line critical profiles\\
Line and smooth conic, transverse & Completely excluded by the same corollary\\
Line and smooth conic, tangent & Tangency, degenerate finite blocks and escaping profiles\\
Three nonconcurrent lines & Completely excluded by the same corollary\\
Three concurrent distinct lines & Constant and nonconstant transverse linear coefficient\\
Double line and a distinct line & Isotropic and nonisotropic repeated-line directions\\
Triple line & Isotropic and nonisotropic repeated-line directions\\
\bottomrule
\end{tabular}
\end{center}
The list is exhaustive by factorization of a plane cubic over $\C$.
The smooth row is now proved here. The ordinary-node gate and its
common residual normal form are also proved. Proposition~\ref{prop:ordinaryfinite}
excludes every finite-critical no-frame member of that family; the
conic, hyperbola and center-axis cases are all excluded by the
preceding propositions. Thus every ordinary-node row is proved here. The other rows remain
proof obligations. For isolated leading singular points, a null
point with $F_2=0$ goes directly to the affine-linear-coordinate table;
otherwise the lower Gaussian quadratic and its cross terms must be
retained in the indicated exclusion. In the repeated-line cases the
singular locus is a whole projective line, and the frame test must be
applied to all its null directions, not to an arbitrarily selected point.

Thus the routing is now explicit. Completing frame extraction still
requires the remaining five geometric families, each with its full
rank or homology-completeness argument. The full affine-linear proof above does not prove those geometric
exclusions. Finite algebra checks do not replace their missing
rank and homology arguments.
